\documentclass[output=paper]{langscibook}
\ChapterDOI{10.5281/zenodo.22880242}
\title{\textit{O baneo visoas} – The young breadfruit tree}
\author{Ruth Sima Rigamu and Joyce Maion}
\abstract{}
\begin{document}
\maketitle\label{chap:8}
\begin{figure}[H]
    \includegraphics[trim=0 50 0 50,clip,width=.6\linewidth]{figures/8_1.jpg}
    \caption{\textit{Baneo visoas} -- Young breadfruit tree}
    \label{fig:8_1}
\end{figure}
\begin{xltabular}{\textwidth}{lXX}
    (\ref{ex:text4-1})&\textit{O baneo vai o visoasi, o baneo vai o kare pura, to mei nana tea manoto}&A young breadfruit tree is a breadfruit tree that has only recently started growing and is not mature yet.\\
    (\ref{ex:text4-2})&\textit{\textbf{Mosi naono baneo}}&\textbf{Cutting breadfruit trees}\\
    (\ref{ex:text4-3})&\textit{Bea abana tau paku bona kave goroto,}&When people are about to make a turtle net,\\
    (\ref{ex:text4-4})&\textit{eori repaa paku bona vahiihio.}&they have a meeting.\\
    (\ref{ex:text4-5})&\textit{E beera repaa sue,}&Then the chief says,\\
    (\ref{ex:text4-6})&\textit{``Bara tobara visoasi, eara kahi paku rara a kave goroto.}&``Well, young boys, we are going to make a turtle net.\\
    (\ref{ex:text4-7})&\textit{Eam toro nao tea mosi ta maa naono baneo vai ta maa rutarutaa.''}&You must go and cut some small breadfruit trees.''\\
    (\ref{ex:text4-8})&\textit{Eori repaa nao.}&Then they go.\\
\end{xltabular}
\begin{figure}[H]
    \captionsetup[subfigure]{oneside,margin={0cm,0cm}}
    \captionsetup{justification=centering}
    \subfloat[]{\includegraphics[trim=0 0 0 50,clip,width=.48\linewidth]{figures/8_2.jpg}\label{fig:8_2}}\hfil
    \subfloat[]{\includegraphics[trim=0 0 0 0,clip,width=.52\linewidth]{figures/8_3.jpg}\label{fig:8_3}}
    \caption{(a) \textit{Kapa baneo} -- Cover of the breadfruit tree,\\ (b) \textit{Sui} -- Ark Clam, used to scrape the fibres of the breadfruit tree}
\end{figure}
\begin{xltabular}{\textwidth}{lXX}
    (\ref{ex:text4-9})&\textit{\textbf{Hihi bona kapa baneo}}&\textbf{Pulling off the cover of breadfruit trees.}\\
    (\ref{ex:text4-10})&\textit{Beori mosi vakavara maa bona naono baneo vai a rutarutaa,}&When they have finished cutting the small breadfruit trees,\\
    (\ref{ex:text4-11})&\textit{eori repaa dee voosu maa bari.}&they carry them home.\\
    (\ref{ex:text4-12})&\textit{A naono baneo vai a rutarutaa na mataa nana, suku bona na kobusu nana,}&The small breadfruit trees are good, because they are moist,\\
    (\ref{ex:text4-13})&\textit{me na antee nana tea tahihi vamataa.}&and (the fibres) can easily be peeled off.\\
    (\ref{ex:text4-14})&\textit{Beori dee voosu maa bari,}&When they have brought them home, \\
    (\ref{ex:text4-15})&\textit{eori repaa taneo tea kahu bono kapa vaa toonuna,}&they start scraping off the top cover (i.e. the bark),\\
    (\ref{ex:text4-16})&\textit{beori varoava batari bari.}&while they are putting them (the trees) in the sun.\\
    (\ref{ex:text4-17})&\textit{Beve raka, eori repaa hihi pete bari.}&When they are dry, they pull off the fibres.\\
    (\ref{ex:text4-18})&\textit{Eori repaa varoava vaha bana bari, eve re kiksi.}&Then they will put them in the sun again, so that they become strong.\\
    (\ref{ex:text4-19})&\textit{Beve mahaka vakavara,}&When they are completely dry,\\
    (\ref{ex:text4-20})&\textit{eori repaa kahu mana bono kara baneo bona kehaa}&they eventually scrape the breadfruit tree fibres with a shell\\
    (\ref{ex:text4-21})&\textit{to dao riori bona sui.}&that they call \textit{sui} (ark clam).\\
    (\ref{ex:text4-22})&\textit{Murinae eori repaa sirisiri varutaa bono kara baneo,}&After that they tear the breadfruit fibres into small strips,\\
    (\ref{ex:text4-23})&\textit{matatopo vaa tea vinugu.}&in preparation for twisting (strings).\\
\end{xltabular}
\begin{xltabular}{\textwidth}{lXX}
    (\ref{ex:text4-24})&\textit{\textbf{Kave goroto}}&\textbf{Turtle net}\\
    (\ref{ex:text4-25})&\textit{Beori vinugu vakavara bona kanono,}&After they finish twisting the string\\
    (\ref{ex:text4-26})&\textit{eori repaa taneo tea huha bona kave goroto.}&they start tying the turtle net.\footnote{See Chapter \ref{chap:27}.}\\
    (\ref{ex:text4-27})&\textit{Kavara na aba vai to huha vatotoena bata maari,}&Everybody who is involved in tying (the net)\\
    (\ref{ex:text4-28})&\textit{na vaa'antee bata maari bono peho mata kave.}&supplies one mesh of the net.\\
    (\ref{ex:text4-29})&\textit{Evei o vahutate no baneo visoasi.}&So, this is the story of the young breadfruit tree.\\
\end{xltabular}
\end{document}